\documentclass[a4paper,12pt]{article}
\usepackage{bbding,combelow,textcomp,amsfonts,amsthm,amsmath,amssymb,graphicx,color,hyperref,etoolbox}
\usepackage[utf8x]{inputenc}
\usepackage[lined,boxed,commentsnumbered]{algorithm2e}
\usepackage{mathtools}
\newtoggle{story}


%%%%%%%%%%%%%%%%%%%%%%%%%%%%%%%%%%%%%%%%%%
%% Customisation options --- update as necessary
%%%%%%%%%%%%%%%%%%%%%%%%%%%%%%%%%%%%%%%%%%

%% Course details: title, semester, instructors

\newcommand{\courseTitle}{Introduction to Probability Theory (Math 2502)}
\newcommand{\semester}{Spring 2026}
\newcommand{\instructorA}{Shagnik Das}
\newcommand{\instructorB}{}

%% Sheet number

\newcommand{\sheetNumber}{2}


%% Submission details

\newcommand{\dueDate}{13:00, Mar 26th.}
\newcommand{\lateWarning}{Submissions that are late will receive no credit.}


%% Display irrelevant footnotes?  \toggletrue for yes, \togglefalse for no

\toggletrue{story}

%%%%%%%%%%%%%%%%%%%%%%%%%%%%%%%%%%%%%%%%%%
%% End of customisation options
%%%%%%%%%%%%%%%%%%%%%%%%%%%%%%%%%%%%%%%%%%

\pdfpagewidth 8.5in
\pdfpageheight 11in
\topmargin -1in
\headheight 0in
\headsep 0in
\textheight 8.5in
\textwidth 6.5in
\oddsidemargin 0in
\evensidemargin 0in
\headheight 77pt
\headsep 0in
\footskip .75in

\makeatletter
\newcommand{\buildtitle}[4]{
\begin{flushleft}
{\large
#1
\hfill{}
#2
\par
#3
}
\end{flushleft}
\vskip 4pt
\begin{center}
{\large\bfseries#4\par}
\end{center}
\bigskip
}
\makeatother

\renewcommand{\thesection}{\normalsize\arabic{section}}

\renewcommand{\deg}{\textrm{deg}}
\newcommand{\eps}{\varepsilon}
\newcommand{\ceil}[1]{\left \lceil #1 \right \rceil}

\newcommand{\hint}{\begin{flushright} [Hint at \url{\hintURL}.] \end{flushright}}

\newcommand{\card}[1]{\left| #1 \right|}
\newcommand{\mb}[1]{\mathbb{#1}}
\newcommand{\mc}[1]{\mathcal{#1}}


\newcommand{\story}[1]{\iftoggle{story}{\footnote{#1}}{}}
\newcommand{\storymark}[1][42]{\iftoggle{story}{\footnotemark[#1]}{}}
\newcommand{\storytext}[2][42]{\iftoggle{story}{\footnotetext[#1]{#2}}{}}

\newcommand{\bonus}[2]{\paragraph{Bonus (#1 pt\ifstrequal{#1}{1}{}{s})} #2}


\begin{document}


\buildtitle{\courseTitle{}}{\semester}{\instructorA{} \hfill \instructorB{}}{Exercise Sheet \sheetNumber{}}

\vspace{-0.2in}

\begin{center}
{B13902024 Chang, Yi-Kuei}

{\bf Due date: \dueDate \\ \lateWarning}

\end{center}

\noindent You should try to solve all of the exercises below --- each problem is worth $10$ points. Make sure to clearly justify your answers, defining everything precisely and stating any results you are using. You should then submit your solutions in Gradescope on NTU COOL, either as a typed PDF or a clear scan of your handwritten answers, and should indicate where in the file the solutions to each exercise are. If you have difficulties with Gradescope, please send your solutions by e-mail.

\paragraph{Exercise 1}  Let $E_1, E_2, \hdots, E_n$ be mutually independent events in a probability space.
\begin{itemize}
	\item[(a)] Show that for any choice of $F_i \in \{ E_i, E_i^c \}$, $1 \le i \le n$, the events $F_1, F_2, \hdots, F_n$ are mutually independent.
	\item[(b)] Deduce that if $F$ is an event formed from $E_1, E_2, \hdots, E_{n-1}$ by taking complements, unions, and intersections,\story{For example, we could have $F = \left( E_1 \cap E_3^c \cap E_4 \right) \cup \left(E_1^c \cap \left( E_2 \cup E_4^c \cup E_6 \right) \right)$.} then $F$ is independent of $E_n$.
\end{itemize}

\paragraph{Solution:} 
\begin{itemize}
	\item [(a)]	Suppose for \(1 \le i \le k\), \(F_i = E_i\) and for \(k + 1 \le i \le n\), \(F_i = E_i^C\). If this not holds, then we can re-label \(\left\{ F_i \right\} \) to make it true. Now we prove on induction that for any \(n \ge 2\), if \(E_1, \dots , E_n\) are mutually independent, and \(F_i = E_i\) for all \(1 \le i \le k\) and \(F_i = E_i^c\) for all \(k + 1 \le i \le n\) for some \(k \le n\), then \(F_1, F_2, \dots , F_n\) are mutually independent. 
	\begin{itemize}
		\item Base case: For \(n = 2\), since \(E_1, E_2\) are mutually independent, so
		\[
			\mathbb{P} (E_1 \cap E_2) = \mathbb{P} (E_1) \mathbb{P} (E_2).
		\]
		Also, 
		\begin{align*}
			\mathbb{P} (E_1 \cap E_2^c) &= \mathbb{P} (E_1) - \mathbb{P} (E_1 \cap E_2) = \mathbb{P} (E_1) - \mathbb{P} (E_1) \mathbb{P} (E_2) = \mathbb{P} (E_1) (1 - \mathbb{P} (E_2)) \\
			&= \mathbb{P} (E_1) \mathbb{P} (E_2^c). \\
			\mathbb{P} (E_1^c \cap E_2^c) &= \mathbb{P} (E_1 \cup E_2)^c = 1 - \mathbb{P} (E_1 \cup E_2) = 1 - (\mathbb{P} (E_1) + \mathbb{P} (E_2) - \mathbb{P} (E_1 \cap E_2)) \\
			&= (1 - \mathbb{P} (E_1)) - \mathbb{P} (E_2) + \mathbb{P} (E_1 \cap E_2) \\
			&= \mathbb{P} (E_1^c) - \mathbb{P} (E_2) + \mathbb{P} (E_1) \mathbb{P} (E_2) = \mathbb{P} (E_1^c) - \mathbb{P} (E_2) (1 - \mathbb{P} (E_1)) \\
			&= \mathbb{P} (E_1^c) - \mathbb{P} (E_2) \mathbb{P} (E_1^c) = (1 - \mathbb{P} (E_2))\mathbb{P} (E_1^c) = \mathbb{P} (E_2^c) \mathbb{P} (E_1^c).
		\end{align*}
		Thus, for \(n = 2\) the statement is true. 
		
		\item Induction step: Now we know 
		\begin{align*}
			\mathbb{P} \left( \bigcap_{i=1}^{n} F_i  \right) &= \mathbb{P} \left( \bigcap_{i=1}^{k} E_i \cap \bigcap_{i=k+1}^{n} E_i^c   \right) = \mathbb{P} \left( \bigcap_{i=1}^{k} E_i \cap \left( \bigcup_{i=k+1}^{n} E_i \right)^c   \right) \\
			&= \mathbb{P} \left( \bigcap_{i=1}^{k} E_i \right) - \mathbb{P} \left( \bigcap_{i=1}^{k} E_i \cap \bigcup_{i=k+1}^{n} E_i \right).   
		\end{align*}
		Define \(C_j = \bigcap_{i=1}^{k} E_i \cap E_j \), and let \(S = \left\{ z \in \mathbb{N} : k + 1 \le z \le n \right\} \), so 
		\begin{align*}
			\mathbb{P} \left( \bigcap_{i=1}^{k} E_i \cap \bigcup_{i=k+1}^{n} E_i \right) &= \mathbb{P} \left( \bigcup_{i=k+1}^{n} C_i \right) = \sum_{\varnothing \neq I \subseteq S} (-1)^{\vert I \vert + 1} \mathbb{P} \left( \bigcap_{i \in I} C_i  \right) \\
			&= \sum_{\varnothing \neq I \subseteq S} (-1)^{\vert I \vert + 1} \mathbb{P} \left( \bigcap_{i=1}^{k} E_i \cap \bigcap_{j \in I} E_j   \right) \\
			&= \sum_{\varnothing \neq I \subseteq S} (-1)^{\vert I \vert + 1} \mathbb{P} \left( \bigcap_{i=1}^k E_i  \right) \mathbb{P} \left( \bigcap_{j \in I} E_j  \right) \\
			&= \mathbb{P} \left( \bigcap_{i=1}^k E_i  \right) \left( \sum_{\varnothing \neq I \subseteq S} (-1)^{\vert I \vert + 1} \mathbb{P} \left( \bigcap_{j \in I} E_j  \right) \right) \\
			&= \mathbb{P} \left( \bigcap_{i=1}^k E_i  \right) \mathbb{P} \left( \bigcup_{i=k+1}^{n} E_i  \right).          
		\end{align*} 
		This gives 
		\begin{align*}
			\mathbb{P} \left( \bigcap_{i=1}^{n} F_i  \right) &= \mathbb{P} \left( \bigcap_{i=1}^{k} E_i \right) - \mathbb{P} \left( \bigcap_{i=1}^{k} E_i \cap \bigcup_{i=k+1}^{n} E_i \right) \\
			&= \mathbb{P} \left( \bigcap_{i=1}^{k} E_i  \right) - \mathbb{P} \left( \bigcap_{i=1}^{k} E_i  \right) \mathbb{P} \left( \bigcup_{i=k+1}^{n} E_i  \right) \\
			&= \mathbb{P} \left( \bigcap_{i=1}^{k} E_i  \right)  \left( 1 - \mathbb{P} \left( \bigcup_{i=k+1}^{n} E_i  \right)  \right) \\
			&= \mathbb{P} \left( \bigcap_{i=1}^{k} E_i  \right)  \mathbb{P} \left( \bigcup_{i=k+1}^{n} E_i  \right)^c = \mathbb{P} \left( \bigcap_{i=1}^{k} E_i  \right)  \mathbb{P} \left( \bigcap_{i=k+1}^{n} E_i^c  \right),      
		\end{align*}
		and by induction hypothesis,
		\[
			\mathbb{P} \left( \bigcap_{i=1}^{n} F_i  \right) = \mathbb{P} \left( \bigcap_{i=1}^{k} E_i  \right)  \mathbb{P} \left( \bigcap_{i=k+1}^{n} E_i^c  \right) = \prod _{i=1}^k \mathbb{P} (E_i) \prod _{i = k + 1}^n \mathbb{P} \left( E_i^c \right) = \prod _{i=1}^n \mathbb{P} (F_i), 
		\]
		so we're done.
	\end{itemize}
	\item [(b)] We know
	\[
		F = \bigcup_{i=1}^{k} A_i, \text{ where } A_i = \bigcap_{j=1}^{s_i} F_{a_{i, j}} \text{ s.t. } F_{a_{i, j}} \in \left\{ E_{a_{i, j}}, E_{a_{i, j}}^c \right\} \text{ and } 1 \le a_{i, j} \le n-1 \text{ for all } i, j.      
	\]
	Hence,
	\begin{align*}
		\mathbb{P} (F \cap E_n) &= \mathbb{P} \left( \left( \bigcup_{i=1}^{k} A_i  \right) \cap E_n  \right) = \mathbb{P} \left( \bigcup_{i=1}^{k} (A_i \cap E_n)  \right) \\
		&= \sum_{\varnothing \neq I \subseteq [k]} (-1)^{\vert I \vert + 1} \mathbb{P} \left( \bigcap_{i \in I} (A_i \cap E_n)  \right) = \sum_{\varnothing \neq I \subseteq [k]} (-1)^{\vert I \vert + 1} \mathbb{P} \left( \bigcap_{i \in I} A_i \cap E_n  \right) \\
		&= \sum_{\varnothing \neq I \subseteq [k]} (-1)^{\vert I \vert + 1} \mathbb{P} \left( \bigcap_{i \in I} A_i  \right) \mathbb{P} (E_n) = \mathbb{P} \left( \bigcup_{i=1}^k A_i  \right)  \mathbb{P} (E_n) = \mathbb{P} (F) \mathbb{P} (E_n)   
	\end{align*}
	by (a), and we're done.
\end{itemize}

\paragraph{Exercise 2}  A local TV station hosts a dumpling tasting contest --- in each round, the contestants are served three dumplings --- a pork and cabbage dumpling, a pork and leek dumpling, and a beef dumpling --- in a uniformly random order. They have to identify the flavour of each dumpling, and whoever gets the most dumplings correct in each round progresses to the next round.

After several rounds of fierce competition, three people are left for the Grand Final --- Ho Keun, David, and Mark. However, after eating so many dumplings, their taste buds are completely exhausted, so they cannot taste anything. Instead, they come up with different strategies.

Ho Keun decides to guess that each dumpling is a pork and cabbage dumpling, since that is his favourite flavour. David, on the other hand, guesses a uniformly random flavour for each dumpling, independently of his previous guesses. Mark, though, realises that the flavours cannot be repeated, so instead chooses a uniformly random order of the three flavours for his guesses.

In order to win the Grand Final, a player has to score strictly more points than the other two --- ties mean there will be no winner.

\begin{itemize}
	\item[(a)] Before doing any calculations, which player do you think has the greatest probability of winning?
	\item[(b)] Calculate the probability of each player winning the Grand Final.
\end{itemize}

\paragraph{Solution:}
\begin{itemize}
	\item [(a)] I think Ho Keun will win because his name is more asian, so he knows more about dumplings. 
	\item [(b)] Suppose \(P_1 = \) Ho Kuen, \(P_2 = \) David, \(P_3 = \) Mark, and \(n(P_i)\) means the number of dunmplings that \(P_i\) guesses correctly among three dumplings. Note that \(n(P_1) = 1\) in all cases since one of the three dumplings must be pork and cabbage dumpling, and the other two must not be. Hence, the only way that \(P_1\) win is \(n(P_2) = n(P_3) = 0\). Thus,
	\begin{align*}
		\mathbb{P} (P_1 \text{ win} ) &= \mathbb{P} (n(P_2) = n(P_3) = 0) = \mathbb{P} (n(P_2) = 0) \cdot \mathbb{P} (n(P_3) = 0) \\
		&= \left( \frac{2}{3} \right)^3 \cdot \frac{2}{3!} = \frac{8}{81}, 
	\end{align*}
	since the guess result of \(P_2\) and \(P_3\) are independent. 
	
	Now we compute \(\mathbb{P} (P_2 \text{ win} )\). \(P_2\) wins if and only if
	\[
		n(P_2) \ge 2 \text{ and } n(P_3) < n(P_2) 
	\] 
	since \(n(P_1) = 1\). Thus, 
	\begin{align*}
		\mathbb{P} (P_2 \text{ win}) &= \mathbb{P} (n(P_2) = 3 \cap n(P_3) \in \left\{ 0, 1 \right\} ) + \mathbb{P} (n(P_2) = 3 \cap n(P_3) \in \left\{ 0, 1, 2 \right\} ) \\
		&= \mathbb{P} (n(P_2) = 2) \left( \mathbb{P} (n(P_3) = 0) + \mathbb{P} (n(P_3) = 1) \right) + \mathbb{P} (n(P_3) = 3) (1 - \mathbb{P} (n(P_3) = 3)) \\
		&= \binom{3}{2} \left( \frac{1}{3} \right)^2 \left( \frac{2}{3} \right) \cdot \left( \frac{2}{3!} + \frac{3}{3!} \right) + \left( \frac{1}{3} \right)^3 \left( 1 - \frac{1}{3!} \right) = \frac{35}{162}.      
	\end{align*} 

	Now we compute \(\mathbb{P} (P_3 \text{ win} )\). Note that \(P_3\) wins if and only if 
	\[
		n(P_3) \ge 2 \text{ and }  n(P_2) < n(P_3).
	\]     
	However, note that \(\mathbb{P} (n(P_3) = 2) = 0\) since for each order of three dumplings, it is only possible that Mark guess right \(0, 1, 3\) dumplings. Thus,
	\begin{align*}
		\mathbb{P} (P_3 \text{ win} ) &=  \mathbb{P} (n(P_3) = 3 \cap n(P_2) \in \left\{ 0, 1, 2 \right\} ) = \mathbb{P} (n(P_3) = 3) \left( 1 - \mathbb{P} (n(P_2) = 3) \right) \\
		&= \frac{1}{3!} \left( 1 - \left( \frac{1}{3} \right)^3  \right) = \frac{13}{81}. 
	\end{align*}  

\end{itemize}

\paragraph{Exercise 3}  Recall the unfortunate gambler from the example in lectures. After his latest bankruptcy, he is forced to move to a less reputable part of town, where the local casino rigs the games. This time, for every round he plays, with probability $p$, where $0 < p < \frac12$, he wins, in which case he gains \$1, while with probability $1 - p$, he loses, in which case he loses \$1. As before, the results of each game are mutually independent.

Suppose he enters with \$k in his pocket, and will keep playing until either he runs out of money, or he has a total of \$n, where $0 < k < n$. What is the probability that he will run out of money?

\paragraph{Solution: } Suppose \(q_i\) is the probability of running out of money starting from \(\$ i\). Then, \(q_1 = 0\) and \(q_n = 0 \). Now suppose \(E_1 = \) next round win and \(E_2 = \) next round lose, then \(S = E_1 \cup E_2 \), ans suppose \(F = \text{ running out of money starting at } \ell\), then by law of total probability:
\[
	\mathbb{P} (F) = \mathbb{P} (E_1) \mathbb{P} (F \mid E_1) + \mathbb{P} (E_2) \mathbb{P} (F \mid E_2),
\]
which is equivalent to
\[
	q_{\ell } = p q_{\ell + 1} + (1 - p) q_{\ell - 1} \text{ for all } 1 \le \ell \le n - 1. 
\]  
We want to compute \(q_k\). Hence, we have 
\[
	\begin{dcases}
		q_{\ell + 1} = \frac{1}{p} q_{\ell } + \frac{p-1}{p} q_{\ell - 1} &\text{for } 1 \le \ell \le n - 1 \\
		q_0 = 0 \\
		q_n = 1.  
	\end{dcases}
\]   
and thus the characteristic function of this recurrence relation is 
\[
	x^2 = \frac{1}{p} x + \frac{p-1}{p},
\]
which has roots \(1, \frac{1-p}{p}\). Thus,
\[
	q_{\ell } = c + d \left( \frac{1-p}{p} \right)^{\ell }  
\] for some constant \(c, d\). Thus, by plug \(0, n\) in \(\ell \) we have 
\[
	\begin{dcases}
		1 = c + d \\
		0 = c + d \left( \frac{1-p}{p} \right)^n.
	\end{dcases}
\] 
Thus, 
\[
	c = 1 - \frac{1}{1 - \left( \frac{1-p}{p} \right)^n }, \quad d = \frac{1}{1 - \left( \frac{1-p}{p} \right)^n }.
\]
Thus,
\begin{align*}
	q_k &= 1 - \frac{1}{1 - \left( \frac{1-p}{p} \right)^n } + \frac{1}{1 - \left( \frac{1-p}{p} \right)^n } \left( \frac{1-p}{p} \right)^k = \frac{\left( \frac{1-p}{p} \right)^k - \left( \frac{1-p}{p} \right)^n  }{1 - \left( \frac{1-p}{p} \right)^n }. 
\end{align*}

\paragraph{Exercise 4}  Consider the experiment of choosing a point $z$ from the interval $(0,1]$ uniformly at random. In this setting, we have the (infinite) sample space $S = (0,1]$. Given any interval $(x,y] \subseteq S$, we can define the event of $z$ belonging to the interval; that is, $E_{x,y} = \{ z \in (x,y] \} = \{ x < z \le y \}$. Under the uniform distribution, the probability of this event is $\mathbb{P}\left( E_{x,y} \right) = y - x$.

More generally (as required by the third axiom of probability), if there is a disjoint union of intervals $(x_1, y_1] \cup (x_2, y_2] \cup \hdots$ with $x_1 \le y_1 < x_2 \le y_2 < \hdots$, then we can define the event $E$ that $z$ lies in the union of these intervals, and $\mathbb{P}(E) = \sum_{i=1}^\infty \mathbb{P}(E_{x_i,y_i}) = \sum_{i=1}^\infty y_i - x_i$.

Show that for any values $p_1, p_2, p_3, \hdots$ satisfying $0 \le p_i \le 1$ and $\sum_{i=1}^{\infty} p_i = \infty$, there are events $E_1, E_2, E_3, \hdots$ with $\mathbb{P}(E_i) = p_i$ and $\mathbb{P}( \limsup_{n \to \infty} E_n ) = 1$. Recall that $\limsup_{n \to \infty} E_n$ is the event that infinitely many of the events $E_i$ occur.

\paragraph{Solution: } 
We define \(F(s) = s - \lfloor s \rfloor\), and then we define 
\[
	L: \mathbb{R} ^2 \to \left\{ S: S \subseteq [0,1)  \right\}
\]
by 
\[
	L(a, b) \coloneqq \begin{dcases}
		[F(a), F(b)), &\text{ if } F(a) \leq F(b) ;\\
		[F(a), 1) \cup [0, F(b)], &\text{ if } F(a) > F(b).
	\end{dcases}
\]

Now suppose the sample space \(S = [0, 1)\) and the probability is the uniform distribution, and \(p_0 = 0\), then we can define event \(E_i\) by  
\[
	E_i = \left\{ x: x \in L \left( \sum_{k=0}^{i-1} p_k, \sum_{k=0}^i p_k   \right)  \right\} \text{ for all } i \in \mathbb{N}. 
\] 
Hence, we can check that 
\[
	\mathbb{P} (E_i) = p_i \text{ for all } i \in \mathbb{N}.  
\]
This can be thought by letting \(I_i = \left( \sum_{k=0}^{i-1} p_k, \sum_{k=0}^i p_k   \right] \), then the length of \(I_i\) is \(p_i \le 1\), and \(L \left( \sum_{k=0}^{i-1} p_k, \sum_{k=0}^i p_k   \right)\) is \(I_i\) modulo \(1\), which is the set of all decimal part \(I_i\) contains. Now since \(\sum_{i=1}^{\infty} p_i = \infty \), so 
\[
	\bigcup_{i=1}^{\infty } I_i = \left\{ x \in \mathbb{R} : x > 0 \right\}.  
\]      
Now for any \(x \in [0, 1)\), we know for all \(k \in \mathbb{N} \cup \left\{ 0 \right\} \), \(x + k \in I_j\) for a unique \(j \in \mathbb{N} \) (since \(I_p, I_q\) are disjoint for distinct \(p, q\)). Also, for any distinct \(k_1, k_2 \in \mathbb{N} \cup \left\{ 0 \right\}  \),  since
\[
	\vert (x + k_1) - (x + k_2) \vert = \vert k_1 - k_2 \vert \ge 1,  
\]   
so if \(x + k_1 \in I_{j_1}\) and \(x + k_2 \in I_{j_2}\), then \(j_1 \neq j_2\), so there exists infinitely many \(j \in \mathbb{N} \) s.t. \(x + k \in I_j\) for some \(k \in \mathbb{N} \cup \left\{ 0 \right\} \). Also, for \(x \in [0, 1)\), \(x + k \in I_j\) for some \(k \in \mathbb{N} \cup \left\{ 0 \right\} \) iff
\[
	x \in L \left( \sum_{r=0}^{j-1} p_r, \sum_{r=0}^j p_r   \right), 
\]    
so we have shown that for any \(x \in [0, 1)\), there are infinitely many \(j \in \mathbb{N} \) has 
\[
	x \in L \left( \sum_{r=0}^{j-1} p_r, \sum_{r=0}^j p_r   \right),
\]   
i.e. there are infinitely many \(j \in \mathbb{N} \) has \(x \in E_j\). Thus, 
\[
	\mathbb{P} (\limsup_{n \to \infty} E_n ) = 1.
\]  

%\section*{Hints}  The following QR codes contain hints for some of the homework exercises.  You should be able to decode them using any QR code scanner, including this one: \url{https://zxing.org/w/decode.jspx}.

%\paragraph{Exercise 4} Also available at \url{https://i.ibb.co/MDgMvpXB/S02E4.png}.

%\begin{center}
	%\includegraphics[scale=0.5]{Hints/S02E4.png}
%\end{center}

\end{document}